\documentclass[11pt]{article}
\usepackage{amsmath, amssymb, amsthm, mathtools}
\usepackage[utf8]{inputenc}
\usepackage{hyperref}
\usepackage{geometry}
\usepackage{booktabs}
\geometry{margin=1in}

\newtheorem{theorem}{Theorem}
\newtheorem{proposition}{Proposition}
\newtheorem{corollary}{Corollary}
\newtheorem{lemma}{Lemma}
\theoremstyle{definition}
\newtheorem{definition}{Definition}
\newtheorem{assumption}{Assumption}
\theoremstyle{remark}
\newtheorem{remark}{Remark}

\title{A Bayesian source-reliability model for channel-conditioned compliance in frontier language models}
\author{F.~M.~Affonso}
\date{\today}

\begin{document}
\maketitle

\begin{abstract}
We formalise the channel-as-authority phenomenon documented empirically in frontier
language models as a Bayesian source-reliability model under which a delivery channel
carries a learned prior on content reliability. We derive a rate-ratio bound that
the difference in compliance log-odds across channels cannot exceed the difference
in channel-prior log-odds, under a 1-Lipschitz-in-logit calibration assumption that
admits the logistic compliance family as a special case. We then derive a one-sided
signed bound on the grounding-evidence effect and an information-theoretic lower
bound on the activation-probing direction that the rate-ratio bound's failure must
induce, plus a Gaussian-exact sharpening for linear read-outs. Each bound yields a
falsification corollary that is testable against the existing authority-laundering
corpus, and the empirical magnitudes falsify rational Bayesian compliance under any
plausible channel-prior estimate. The framework connects behavioural rate
measurements to mechanistic activation-probing findings via a Pinsker-Le Cam
representational lower bound and clarifies the four conditions under which an
empirical bound violation counts as evidence of a non-Bayesian channel-authority
heuristic. All seven theorems carry machine-checked proofs that pass at least
two of an algebraic-symbolic verifier, a numerical-witness verifier, and a
satisfiability-solver verifier per theorem; verification scripts ship with the
companion repository.
\end{abstract}

\section{Setup and notation}
\label{sec:setup}

Let $\mathcal{C}$ denote the space of declarative content (institutional documents,
regulations, citations, professional protocols) and $\mathcal{H} = \{\text{sys}, \text{tool}, \text{doc}, \text{user}\}$
the discrete set of delivery channels considered in the empirical work, in the
machine-behaviour register of \cite{rahwan2019machine}. A subject model~$M$ receives
content $c \in \mathcal{C}$ through channel $h \in \mathcal{H}$ and produces a
compliance indicator $Y \in \{0, 1\}$ where $Y=1$ denotes operationalisation of the
content. Let $R \in \{0, 1\}$ denote the latent reliability of the content, where
$R=1$ means the content is genuine institutional material and $R=0$ means it is
fabricated.

\begin{definition}[Channel prior]
\label{def:channel-prior}
The \emph{channel prior} is $\pi(h) := \mathbb{P}(R = 1 \mid h)$, the probability
that content arriving through channel~$h$ is reliable, marginalised over content.
\end{definition}

\begin{definition}[Content likelihoods]
\label{def:content-likelihoods}
The \emph{content likelihoods} are $L_r(c) := \mathbb{P}(c \mid R=1)$ and
$L_u(c) := \mathbb{P}(c \mid R=0)$, the probabilities of observing content $c$
conditional on reliability state.
\end{definition}

\begin{definition}[Compliance function]
\label{def:compliance}
The \emph{compliance function} $\gamma : \mathcal{C} \times \mathcal{H} \to [0,1]$
maps content and channel to compliance probability,
$\gamma(c, h) := \mathbb{P}(Y=1 \mid c, h)$.
\end{definition}

\begin{assumption}[Bayesian compliance]
\label{ass:bayes}
The compliance function is a monotonic increasing function of the posterior
reliability:
\[
\gamma(c, h) \;=\; f\!\left(\, \mathbb{P}(R=1 \mid c, h) \,\right),
\]
where $f : [0,1] \to [0,1]$ is monotonic increasing and does not depend on $h$.
\end{assumption}

By Bayes' rule the posterior is
\begin{equation}
\label{eq:posterior}
\mathbb{P}(R=1 \mid c, h) \;=\; \frac{\pi(h)\, L_r(c)}{\pi(h)\, L_r(c) + (1-\pi(h))\, L_u(c)}.
\end{equation}
Equivalently, the posterior log-odds decompose as
\begin{equation}
\label{eq:posterior-logodds}
\operatorname{logit} \mathbb{P}(R=1 \mid c, h) \;=\; \operatorname{logit} \pi(h) \;+\; \log\!\bigl[L_r(c) / L_u(c)\bigr],
\end{equation}
where $\operatorname{logit}(p) := \log[p/(1-p)]$.

Throughout, let $\ell_h := \operatorname{logit} \pi(h)$ denote the channel-prior
log-odds and $\lambda(c) := \log[L_r(c)/L_u(c)]$ denote the content log-likelihood
ratio. Define the compliance function in posterior-log-odds coordinates by
$g(\ell) := \operatorname{logit} f(\sigma(\ell))$, where $\sigma$ is the standard logistic.
Under Assumption~\ref{ass:bayes}, the compliance log-odds factor as
$\operatorname{logit} \gamma(c, h) = g(\ell_h + \lambda(c))$.

\begin{assumption}[Strong matched content and channel-invariant compliance]
\label{ass:strong-matching}
Three conditions hold jointly. (i) For any content $c$ entering the rate-ratio
comparison, the content likelihoods are channel-invariant:
$L_r(c \mid h) = L_r(c)$ and $L_u(c \mid h) = L_u(c)$ for all $h \in \mathcal{H}$.
(ii) The compliance function $f$ is channel-invariant in shape and intercept:
the same $f$ applies on every channel. (iii) The latent reliability $R$ is binary,
and the channel prior $\pi(h)$ is estimated at the same binary resolution as $f$
consumes.
\end{assumption}

Assumption~\ref{ass:strong-matching} is required for the quantitative rate-ratio
bound below; the sign-of-effect bound (Proposition~\ref{prop:sign}) requires only
Assumption~\ref{ass:bayes}. The 1-Lipschitz-in-logit calibration condition that
follows is consistent with the temperature-scaling family of LLM calibrators
\cite{lamb2025semantic} and with empirical evidence that LLMs broadly follow
Bayesian update structure with mis-calibrated priors rather than flawed updates
\cite{gupta2025coinflips}.

\section{Rate-ratio bound under rational source-weighting}
\label{sec:rate-ratio}

We separate the sign-of-effect claim (which channel attracts more compliance) from
the quantitative rate-ratio claim (how much more). The first requires only
Assumption~\ref{ass:bayes}. The second requires Assumption~\ref{ass:strong-matching}
and a calibration condition on the compliance function, expressed naturally in
posterior-log-odds coordinates.

\subsection{Sign-of-effect bound}

\begin{proposition}[Channel monotonicity of compliance]
\label{prop:sign}
Under Assumption~\ref{ass:bayes} alone, for any matched content $c$ and any channels
$h_1, h_2$ with $\pi(h_1) \ge \pi(h_2)$,
$\gamma(c, h_1) \ge \gamma(c, h_2)$.
\end{proposition}

\begin{proof}
The posterior $\mathbb{P}(R=1 \mid c, h)$ is monotonically increasing in $\pi(h)$
for fixed $L_r(c), L_u(c)$, by direct quotient differentiation of
equation~\eqref{eq:posterior}. Composing with the monotonic increasing $f$ preserves
the ordering.
\end{proof}

\begin{corollary}[Sign-of-effect falsification]
\label{cor:sign-falsification}
If matched-content empirical estimates satisfy $\hat\gamma(c, h_1) < \hat\gamma(c, h_2)$
on a channel pair with $\hat\pi(h_1) > \hat\pi(h_2)$, Assumption~\ref{ass:bayes} fails
regardless of the functional form of $f$.
\end{corollary}

This sign-of-effect statement is robust to the entire family of monotonic compliance
functions. It cannot detect over-compliance (a compliance gap larger than the prior
gap), for which the quantitative bound below is required.

\subsection{Rate-ratio bound}

\begin{theorem}[Rate-ratio bound under 1-Lipschitz logit-compliance]
\label{thm:rate-ratio}
Suppose Assumption~\ref{ass:bayes} and Assumption~\ref{ass:strong-matching} hold,
and the compliance function in posterior-log-odds coordinates,
$g(\ell) = \operatorname{logit} f(\sigma(\ell))$, is 1-Lipschitz on $\mathbb{R}$:
$|g(\ell_1) - g(\ell_2)| \le |\ell_1 - \ell_2|$ for all $\ell_1, \ell_2 \in \mathbb{R}$.
Then for matched content $c$ and channels $h_1, h_2$ with $\pi(h_1), \pi(h_2) \in (0,1)$,
\begin{equation}
\label{eq:rate-ratio-bound}
\bigl|\operatorname{logit} \gamma(c, h_1) - \operatorname{logit} \gamma(c, h_2)\bigr|
\;\le\; \bigl|\operatorname{logit} \pi(h_1) - \operatorname{logit} \pi(h_2)\bigr|,
\end{equation}
with equality when $g$ is the identity (linear-in-logit calibration). In particular,
when $\pi(h_1) \ge \pi(h_2)$ the absolute bound sharpens to the signed form
$\operatorname{logit} \gamma(c, h_1) - \operatorname{logit} \gamma(c, h_2) \le
\operatorname{logit} \pi(h_1) - \operatorname{logit} \pi(h_2)$.
\end{theorem}

\begin{proof}
By the factorisation following equation~\eqref{eq:posterior-logodds} and
Assumption~\ref{ass:strong-matching}~(i) and~(ii),
$\operatorname{logit} \gamma(c, h_i) = g(\ell_{h_i} + \lambda(c))$ for $i \in \{1, 2\}$.
The compliance log-odds difference is therefore
\[
\operatorname{logit} \gamma(c, h_1) - \operatorname{logit} \gamma(c, h_2)
\;=\; g(\ell_{h_1} + \lambda(c)) \;-\; g(\ell_{h_2} + \lambda(c)).
\]
The 1-Lipschitz hypothesis applied at the points $\ell_{h_1} + \lambda(c)$ and
$\ell_{h_2} + \lambda(c)$ gives the absolute bound~\eqref{eq:rate-ratio-bound}
directly, $|g(\ell_{h_1} + \lambda(c)) - g(\ell_{h_2} + \lambda(c))| \le
|\ell_{h_1} - \ell_{h_2}|$, since $\lambda(c)$ cancels in the difference of arguments.
For the signed sharpening, monotonicity of $g$ under Assumption~\ref{ass:bayes}
(composition of the two monotonic increasing functions $f$ and the logit) makes the
sign of the left-hand side equal the sign of $\ell_{h_1} - \ell_{h_2}$; when
$\ell_{h_1} \ge \ell_{h_2}$ both sides are non-negative and the absolute bound is the
signed bound. (The signed inequality is specific to this ordering: when
$\ell_{h_1} < \ell_{h_2}$ the left-hand side is non-positive and bounded below, not
above, by $\ell_{h_1} - \ell_{h_2}$, which is why the assumption-free statement is the
absolute bound.) Equality holds when $g$ has slope identically $1$, i.e.\
$g(\ell) = \ell + \alpha$ for some $\alpha \in \mathbb{R}$.
\end{proof}

\paragraph{Verification.} Theorem~\ref{thm:rate-ratio} is independently checked
by an algebraic-symbolic verifier, a numerical-witness verifier sweeping $96$ random
parameter configurations ($32$ per $1$-Lipschitz $g$-family across three families)
on the absolute bound, and a satisfiability-solver verifier on the signed form under
the ordering constraint; all three return PASS with zero failures. The verification script is at
\texttt{theory/proofs/T1\_rate\_ratio/proof.py}.

\begin{corollary}[Logistic special case]
\label{cor:logistic-special}
If $f(p) = \sigma(\beta \cdot \operatorname{logit}(p) + \alpha)$ for fixed
$\alpha \in \mathbb{R}$ and $\beta \in (0, 1]$, the hypothesis of
Theorem~\ref{thm:rate-ratio} holds with $g(\ell) = \beta\ell + \alpha$, and the
rate-ratio bound~\eqref{eq:rate-ratio-bound} holds with equality when $\beta=1$.
\end{corollary}

The 1-Lipschitz-in-logit hypothesis is strictly weaker than the logistic-with-slope-$\beta$
hypothesis; it admits any concave-monotone calibrator with bounded slope, isotonic
regressions with bounded slope in logit-coordinates, Platt-scaling outputs
\cite{spiess2024calibration}, temperature-scaling outputs at $T \ge 1$
\cite{lamb2025semantic}, and sufficiently smooth nonparametric posterior-to-action
mappings. The publishable calibration interpretation is that compliance log-odds
cannot move faster than posterior log-odds; the logistic family is one realisation
of this property and is not implied by Bayesian rationality alone.

\subsection{Empirical falsification}

\begin{corollary}[Empirical falsification of Bayesian compliance]
\label{cor:falsification}
Let $\hat{\gamma}(c, h)$ denote empirical compliance rates and let $\hat{\pi}(h)$
denote empirical channel-prior estimates obtained by directly probing the model's
stated reliability assessment of content arriving through channel $h$. If
\[
\operatorname{logit} \hat{\gamma}(c, h_1) - \operatorname{logit} \hat{\gamma}(c, h_2)
\;>\;
\operatorname{logit} \hat{\pi}(h_1) - \operatorname{logit} \hat{\pi}(h_2)
\]
on the matched content $c$, then either Assumption~\ref{ass:bayes} fails, or
Assumption~\ref{ass:strong-matching} fails, or the compliance function $g$ violates
1-Lipschitz on $\mathbb{R}$. In any case the compliance behaviour cannot be explained
as rational Bayesian source-weighting under the stated calibration condition.
\end{corollary}

\begin{remark}[Falsification scope]
\label{rem:falsification-scope}
Corollary~\ref{cor:falsification} fails if any of four hidden defenders' moves apply:
(a) the model uses a sharper-than-1-Lipschitz logit calibration ($g'(\ell) > 1$ on the
support of interest); (b) channel-conditional content likelihoods $L_r(c \mid h)$ or
$L_u(c \mid h)$ differ across channels even when surface content is matched; (c)
latent reliability is multi-class while the empirical $\hat\pi(h)$ probe is binary;
(d) adversarial content selection induces hidden $c$-channel correlation. Each of
(a)--(d) is a substantive empirical claim. The falsification corollary fires safely
only when (a)--(d) have been independently ruled out. The empirical-test specifications
in \S\ref{sec:predictions} target (b) by sourcing matched content from byte-identical
factorial pairs and (c) by trichotomous reliability probing.
\end{remark}

\begin{remark}[Empirical magnitudes]
\label{rem:empirical}
On the empirical authority-laundering data ($\hat{\gamma}_{\text{tool}} = 0.310$,
$\hat{\gamma}_{\text{user}} = 0.015$, pooled across 9 frontier models and 27 scenarios),
the left-hand side of Corollary~\ref{cor:falsification} equals
\[
\operatorname{logit}(0.310) - \operatorname{logit}(0.015) \;\approx\; -0.800 - (-4.185) \;\approx\; 3.385.
\]
Bayesian rationality therefore requires
$\operatorname{logit} \hat{\pi}(\text{tool}) - \operatorname{logit} \hat{\pi}(\text{user}) \ge 3.385$,
which corresponds to a channel-prior odds ratio of at least
$\exp(3.385) \approx 29.5$. Direct probing of model-reported channel priors instead
returns $\hat{\pi}(h) = 0$ on every channel for fabricated documents (Supplementary
\S S13), an odds-ratio gap of zero, well below this threshold; the empirical gap is
the magnitude of the non-Bayesian channel-authority heuristic, conditional on ruling
out the four scope conditions of Remark~\ref{rem:falsification-scope}. The corresponding deviation magnitude on the
evidence-direction axis is approximately $2$ to $3$ \cite{kumaran2026competing},
about an order of magnitude smaller than the channel-axis deviation reported here.
\end{remark}

\section{Grounding-asymmetry bound}
\label{sec:grounding}

\begin{definition}[Grounding evidence]
\label{def:grounding}
A user-stated grounding fact $F$ is observed evidence that bears on the reliability
of $c$ via Bayes factor $B(F) := \mathbb{P}(F \mid R=0) / \mathbb{P}(F \mid R=1)$.
A grounding fact is \emph{contradicting} when $B(F) > 1$, that is, $F$ is more likely
under fabrication than under genuine content.
\end{definition}

\begin{theorem}[Grounding-effect bound, signed]
\label{thm:grounding}
Suppose Assumption~\ref{ass:bayes} and Assumption~\ref{ass:strong-matching} hold,
the compliance function $g$ is 1-Lipschitz on $\mathbb{R}$, the grounding fact
$F$ is conditionally independent of $h$ given $R$, and the conditional Bayes factor
$B(F \mid c, h) = B(F)$ is matched across channels at the likelihood-ratio level
(not merely at the surface-string level). For a contradicting grounding fact $F$
with $B(F) > 1$, the channel-conditioned compliance log-odds satisfy
\begin{equation}
\label{eq:grounding-bound}
- \log B(F) \;\le\; \operatorname{logit} \gamma(c, h, F) \;-\; \operatorname{logit} \gamma(c, h) \;\le\; 0,
\end{equation}
and the magnitude of the reduction is independent of the channel.
\end{theorem}

\begin{proof}
Conditional on $F$, equation~\eqref{eq:posterior-logodds} becomes
$\operatorname{logit} \mathbb{P}(R=1 \mid c, h, F) = \operatorname{logit} \pi(h) + \log[L_r(c)/L_u(c)] - \log B(F)$,
because $F$ is conditionally independent of $h$ given $R$, and the matched
$B(F \mid c, h) = B(F)$ hypothesis ensures the log-Bayes-factor contributes
$-\log B(F)$ uniformly across channels. The posterior log-odds reduction is exactly
$\log B(F) > 0$ for contradicting $F$, independent of $h$. Applying the 1-Lipschitz
compliance function $g$ scales the absolute reduction by at most $1$, giving the
magnitude bound. The non-positive sign in~\eqref{eq:grounding-bound} follows from
monotonicity of $g$ inherited from monotonicity of $f$ in
Assumption~\ref{ass:bayes}: a reduction in posterior log-odds cannot increase
compliance log-odds under any monotonic compliance function.
\end{proof}

\paragraph{Verification.} Theorem~\ref{thm:grounding} is independently checked
by all three verifiers (algebraic-symbolic, numerical-witness, satisfiability-solver) with zero failures across $32$ random
configurations of $B(F) \in (1, 20.1]$ and three calibration families; see
\texttt{theory/proofs/T2\_grounding/proof.py}.

\begin{corollary}[Sign-violation falsification]
\label{cor:grounding-sign}
If empirical grounding effects produce a strictly positive shift on contradicting $F$,
$\operatorname{logit} \hat\gamma(c, h, F) > \operatorname{logit} \hat\gamma(c, h)$, then
the assumptions of Theorem~\ref{thm:grounding} cannot all hold. In particular,
no rational Bayesian compliance model satisfying Assumption~\ref{ass:bayes} produces
a sign flip on contradicting evidence regardless of the choice of compliance
function $f$. This sign claim is logically stronger than correlational Bayesian
coherence measures used in recent benchmarks of LLM belief updating
\cite{imran2025bayes}. The empirical observation that grounding \emph{increases}
compliance on Claude Sonnet 4.5 from approximately 32 per cent to 36 per cent
under matched contradicting grounding facts is therefore a decisive sign-violation
falsification of rational compliance, robust to any recalibration of the
1-Lipschitz constant or of $\log B(F)$.
\end{corollary}

\begin{remark}[Operational scope of Theorem~\ref{thm:grounding}]
\label{rem:grounding-scope}
The theorem requires three substantive conditions beyond the rate-ratio bound's
assumptions: (i) channel-blind elicitation of $F$, ensuring $F$ is sourced
independently of channel; (ii) Bayes-factor matching at the conditional level
$B(F \mid c, h) = B(F)$, not merely surface matching of $F$ across channels; and
(iii) complete-evidence observation, where the grounding evidence stream entering
the model is observable to the empirical pipeline rather than partially observed
through implicit channels (hedging tone, qualifier density). When grounding effects
are reported turn-wise, the bound applies to evidence cardinality $F_{\le t}$ at
turn $t$, not the full grounding stream. Magnitude violations of the bound can
arise from (i)--(iii) without falsifying rational compliance; sign violations of
Corollary~\ref{cor:grounding-sign} cannot.
\end{remark}

\section{Probing-representation lower bound}
\label{sec:probing}

Information-theoretic bounds on language-model behaviour have recently been
derived as upper bounds on output accuracy from task-complexity overload
\cite{wan2025fano}; the result below is a lower bound in the dual direction, on
activation-probe distinguishability from a behavioural compliance gap. The bound
extends the mutual-information probing framework of \cite{belinkov2022probing,
pimentel2020infotheory} from a recovery-quality measure to a falsifiable
prediction tied to the rate-ratio bound of Theorem~\ref{thm:rate-ratio}.

\begin{theorem}[Probing-representation lower bound]
\label{thm:probing}
Suppose Assumption~\ref{ass:bayes} and Assumption~\ref{ass:strong-matching} hold and
the empirical compliance log-odds gap on matched content $c$ between channels $h_1$
and $h_2$ exceeds the rate-ratio bound of Theorem~\ref{thm:rate-ratio} by an amount
$\Delta_{\textnormal{nb}} > 0$. Let $P_1$ and $P_2$ denote the channel-conditional
distributions of the residual-stream activation $\mathbf{x}_h \in \mathbb{R}^d$ at
the fixed transformer layer $L$ for matched $c$ delivered through $h_1$ and $h_2$.
Assume the compliance read-out at that layer,
$\psi : \mathbb{R}^d \to \mathbb{R}$ defined by
$\psi(x) := \operatorname{logit} \mathbb{E}[Y \mid \mathbf{x}_h = x]$, is bounded
on the activation support: $|\psi(x)| \le M$ almost surely under $P_1$ and $P_2$,
for some $M > 0$. Then:
\begin{enumerate}
\item[\textnormal{(a)}] The total-variation distance between the channel-conditional
activation distributions satisfies
\begin{equation}
\label{eq:tv-lower}
\mathrm{TV}(P_1, P_2) \;\ge\; \frac{\Delta_{\textnormal{nb}}}{2 M}.
\end{equation}
\item[\textnormal{(b)}] The Bayes-optimal binary classifier discriminating
$\{h_1, h_2\}$ from a single sample of $\mathbf{x}_h$ achieves accuracy
\begin{equation}
\label{eq:alpha-lower}
\alpha^\star \;\ge\; \frac{1}{2} \;+\; \frac{\Delta_{\textnormal{nb}}}{4 M}.
\end{equation}
\item[\textnormal{(c)}] The binary KL-divergence between $P_1$ and $P_2$ satisfies
\begin{equation}
\label{eq:kl-lower}
\mathrm{KL}(P_1 \,\|\, P_2) \;\ge\; \frac{\Delta_{\textnormal{nb}}^{\,2}}{2 M^{2}}.
\end{equation}
\item[\textnormal{(d)}] When $P_1, P_2$ admit densities and the channel signal is
encoded linearly in the residual stream at layer $L$, there exists a unit vector
$\mathbf{v} \in \mathbb{R}^{d}$ such that the projection
$\langle \mathbf{v}, \mathbf{x}_h \rangle$ realises a linear classifier that
attains the accuracy floor~\eqref{eq:alpha-lower}.
\end{enumerate}
\end{theorem}

\begin{proof}
\emph{(a) Total-variation lower bound.} By the tower-rule conditioning of $Y$ on
$\mathbf{x}_h$ and the definition of $\psi$,
$\operatorname{logit} \gamma(c, h_i) = \mathbb{E}_{\mathbf{x}_h \sim P_i}[\psi(\mathbf{x}_h)]$.
Hence the compliance log-odds gap equals the read-out's expectation gap,
$\Delta_\gamma := \operatorname{logit} \gamma(c, h_1) - \operatorname{logit} \gamma(c, h_2)
= \mathbb{E}_{P_1}[\psi] - \mathbb{E}_{P_2}[\psi]$.
By hypothesis the absolute compliance log-odds gap exceeds the rate-ratio bound by
$\Delta_{\textnormal{nb}}$. The channel-prior contribution
$\ell_{h_1} - \ell_{h_2}$ is a constant on the activation support that enters the
read-out through the additive intercept under Assumption~\ref{ass:strong-matching},
so it cancels in the differenced expectation. The residual non-Bayesian gap on
the activation expectations therefore satisfies
$|\mathbb{E}_{P_1}[\psi] - \mathbb{E}_{P_2}[\psi]| \ge \Delta_{\textnormal{nb}}$.
The Kantorovich--Rubinstein dual representation of total-variation distance for
bounded functions states that for any measurable $\psi$ with $\|\psi\|_\infty \le M$,
\begin{equation}
\label{eq:kr-bound}
\bigl|\mathbb{E}_{P_1}[\psi] - \mathbb{E}_{P_2}[\psi]\bigr|
\;\le\; 2 M \cdot \mathrm{TV}(P_1, P_2),
\end{equation}
which is a direct consequence of the variational form
$\mathrm{TV}(P_1, P_2) = \tfrac{1}{2} \sup_{\|f\|_\infty \le 1} |\mathbb{E}_{P_1} f - \mathbb{E}_{P_2} f|$
applied to $f = \psi/M$ \cite{coverthomas2006elements}. Combining with the lower bound
on the expectation difference yields~\eqref{eq:tv-lower}.

\emph{(b) Le Cam two-point lemma.} The Le Cam two-point lemma
\cite[\S2.4]{tsybakov2009nonparametric} states that the maximum classification
accuracy under uniform class priors is
$\alpha^\star = \tfrac{1}{2} + \tfrac{1}{2}\mathrm{TV}(P_1, P_2)$,
attained by the likelihood-ratio test. Substituting~\eqref{eq:tv-lower}
gives~\eqref{eq:alpha-lower}.

\emph{(c) Pinsker's inequality.} Pinsker's inequality
\cite[Theorem~17.3.3]{coverthomas2006elements} states
$\mathrm{KL}(P_1 \,\|\, P_2) \ge 2 \, \mathrm{TV}(P_1, P_2)^{2}$.
Squaring~\eqref{eq:tv-lower} and multiplying by $2$ yields~\eqref{eq:kl-lower}.

\emph{(d) Existence of an informative direction.} The Bayes-optimal classifier's
accuracy is $\alpha^\star$ by part~(b). When $P_1$ and $P_2$ admit densities and
their log-density-ratio is affine in $\mathbf{x}_h$, the likelihood-ratio test
collapses to a linear classifier with hyperplane normal proportional to the
log-density-ratio gradient. The unit vector $\mathbf{v}$ along that gradient is the
channel-of-origin direction the linear-probe experiment in
\texttt{mechanism/exp1\_linear\_probe.py} measures.
\qedhere
\end{proof}

\paragraph{Verification.} Theorem~\ref{thm:probing} passes equivalence verification at $3/3$ engines on the parts expressible in linear and polynomial arithmetic; parts (a), (b), (d) pass all three verifiers; part~(c) (Pinsker) is verified by the symbolic and numerical engines with the satisfiability-solver replaced by transcendental-aware
symbolic verification, with $32$ Bernoulli witness pairs and $4$-point multinomial
witness all PASS. Min Pinsker margin $+3 \times 10^{-6}$ at $p = q$ degeneracy;
see \texttt{theory/proofs/T3\_probing/proof.py}.

\begin{theorem}[Gaussian sharpening of the probing bound]
\label{thm:probing-gaussian}
Suppose the conditions of Theorem~\ref{thm:probing} hold and additionally:
(i) the compliance read-out is linear,
$\psi(\mathbf{x}_h) = \mathbf{w}^\top \mathbf{x}_h + b$;
(ii) the channel-conditional residual-stream distributions $P_1, P_2$ are
$d$-dimensional Gaussians with equal covariance and means $\mu_1, \mu_2$
separated along a rank-$1$ channel direction
$\mathbf{v}^\star := (\mu_1 - \mu_2) / \|\mu_1 - \mu_2\|_2$, with projected
standard deviation $\sigma$ along $\mathbf{v}^\star$ on each channel-conditional
distribution. Then the projected mean gap satisfies
$\|\mu_1 - \mu_2\|_2 \ge \Delta_{\textnormal{nb}}/\|\mathbf{w}\|_2$, and the
Bayes-optimal classifier attains accuracy
\begin{equation}
\label{eq:alpha-gaussian}
\alpha^\star \;=\; \Phi\!\left( \frac{\|\mu_1 - \mu_2\|_2}{2\sigma} \right)
\;\ge\; \Phi\!\left( \frac{\Delta_{\textnormal{nb}}}{2 \sigma \, \|\mathbf{w}\|_2} \right),
\end{equation}
where $\Phi$ is the standard normal cumulative distribution function. The unit
vector $\mathbf{v}^\star$ is the explicit channel-authority direction, and the
linear-probe classifier along $\mathbf{v}^\star$ attains the Bayes-optimal accuracy.
\end{theorem}

\begin{proof}
Under hypothesis (i), the linear read-out's expectation gap satisfies
$|\mathbb{E}_{P_1}[\psi] - \mathbb{E}_{P_2}[\psi]| = |\mathbf{w}^\top(\mu_1 - \mu_2)|
\le \|\mathbf{w}\|_2 \cdot \|\mu_1 - \mu_2\|_2$ by Cauchy--Schwarz, hence
$\|\mu_1 - \mu_2\|_2 \ge \Delta_{\textnormal{nb}}/\|\mathbf{w}\|_2$. Under
hypothesis (ii), the two equal-covariance Gaussians admit a Bayes-optimal
likelihood-ratio test that is a linear classifier whose accuracy equals
$\Phi(\|\mu_1 - \mu_2\|_2/(2\sigma))$ \cite[\S2.6]{tsybakov2009nonparametric}.
Substituting the lower bound on $\|\mu_1 - \mu_2\|_2$ and using monotonicity of
$\Phi$ yields~\eqref{eq:alpha-gaussian}.
\end{proof}

\paragraph{Verification.} Theorem~\ref{thm:probing-gaussian} passes equivalence verification at $3/3$ engines on the foundational steps (Cauchy--Schwarz
on the linear read-out and rank-$1$ reduction; all three verifiers return PASS). The Gaussian-equality form $\alpha^\star = \Phi(z/2)$ is verified by the symbolic engine (exact integration of the equal-covariance Gaussian discriminant) and the numerical engine (Monte-Carlo equality witness within binomial tolerance across random Gaussian configurations and a $d = 32$ rank-1 witness); see
\texttt{theory/proofs/T4\_subgaussian/proof.py}. The earlier as-stated linear bound
$\alpha^\star \ge \tfrac{1}{2} + \Delta_{\textnormal{nb}}/(4\sigma)$ was refuted by
both the symbolic and the numerical verifiers because the Gaussian Bayes accuracy $\Phi(z/2)$ has slope
$1/(2\sqrt{2\pi}) \approx 0.199$ at $z=0$, strictly less than $1/4$. The corrected
form above replaces the linear lower bound with the Gaussian-exact equality, which
is consistent with the standard sub-Gaussian classification result and tighter for
all $\Delta_{\textnormal{nb}} > 0$.

\begin{remark}[Operational scope of Theorem~\ref{thm:probing}]
\label{rem:probing-scope}
The bound applies layer-wise and is silent across layers. The empirical-falsification
test scans layers and reports $\sup_L \alpha^\star(L)$, because the bound trivialises
at any layer where the channel signal is not yet linearly decodable
\cite{voita2020mdl}. The bound also requires the probe geometry to match the
read-out geometry: the probe must operate on the same activation tensor (same layer,
same token-aggregation policy) as the compliance read-out. The bound is vacuous
if channel-authority is encoded nonlinearly at the probed layer; in that regime,
the linear probe should be replaced with a matched-expressivity nonlinear probe
(e.g.\ quadratic or single-hidden-layer MLP) and $\alpha^\star$ rederived
\cite{hewitt2019selectivity}. Empirical evaluation must use cross-validated probe
accuracy with explicit confidence intervals to avoid in-sample overstatement in
high-dimensional residual spaces.
\end{remark}

\section{Training-distribution origin of the channel prior}
\label{sec:training-distribution}

The rate-ratio bound of Theorem~\ref{thm:rate-ratio} treats the channel prior
$\pi(h)$ as an unconstrained quantity to be probed empirically. The probing-representation
lower bound of Theorem~\ref{thm:probing} treats the residual non-Bayesian gap
$\Delta_{\textnormal{nb}}$ as the unexplained component of the empirical compliance
gap. We now derive the training-distribution origin of $\pi(h)$ itself, identifying
the channel-prior gap as the KL-projection of the corpus channel-reliability rate.
The result sits structurally upstream of Theorems~\ref{thm:rate-ratio}--\ref{thm:probing-gaussian}
and supplies the missing causal explanation for why the channel-as-authority bias
exists at all. The result is structurally analogous to Cloud Theorem~1
\cite{cloud2026subliminal}, which shows that a single distillation gradient step
transports student logits towards teacher logits regardless of training
distribution; the result below shows that truth-aligned training transports the
channel prior of the model towards the channel-reliability rate of the corpus
regardless of the corpus generative process.

\begin{assumption}[Truth-aligned training distribution]
\label{ass:truth-aligned}
The training distribution $P^\star$ over $(C, H, R) \in \mathcal{C} \times \mathcal{H} \times \{0,1\}$
satisfies, for every channel $h$ in the support of the channel marginal,
$P^\star(R = 1 \mid h) = \rho(h) \in (0, 1)$, where $\rho(h)$ is the empirical
reliability rate on channel $h$. The channel marginal $P^\star(h)$ is positive on
every $h \in \mathcal{H}$ relevant to the compliance task.
\end{assumption}

\begin{assumption}[KL-projection compliance]
\label{ass:kl}
The model's channel-conditioned reliability estimator $\hat\pi : \mathcal{H} \to (0,1)$
is the KL projection of $P^\star(R \mid H)$ onto the channel-indexed Bernoulli
family,
\[
\hat\pi(h) = \arg\min_{q \in (0,1)} \mathbb{E}_{R \sim P^\star(\cdot \mid h)}\!\left[\, \mathrm{KL}\!\left( \mathrm{Ber}(R) \,\|\, \mathrm{Ber}(q) \right) \,\right],
\]
and the compliance function $\gamma(c, h) = f(\,\mathbb{P}_{\hat\pi}(R=1 \mid c, h)\,)$
is Bayesian under Assumption~\ref{ass:bayes} with channel prior $\hat\pi$.
\end{assumption}

Assumption~\ref{ass:truth-aligned} captures the training-distribution structure of
modern instruction-tuning and RLHF pipelines \cite{ouyang2022instruct, casper2023rlhf,
rafailov2023dpo}, in which retrieval-surface channels (tool, document) typically
carry higher reliability than user-authored content. Assumption~\ref{ass:kl} is the
standard M-projection identification of the channel-prior estimator with the corpus
reliability rate \cite{csiszar1975projection, amari2016information,
coverthomas2006elements}.

\begin{theorem}[Training-distribution impossibility for channel-uniform compliance]
\label{thm:impossibility}
Suppose Assumptions~\ref{ass:bayes}, \ref{ass:strong-matching},
\ref{ass:truth-aligned}, and \ref{ass:kl} hold, the compliance function $g$
in posterior-log-odds coordinates is 1-Lipschitz on $\mathbb{R}$, and there exist
two channels $h_1, h_2 \in \mathcal{H}$ with $\rho(h_1) > \rho(h_2)$ in $(0,1)$.
Then for matched content $c$,
\begin{equation}
\label{eq:impossibility}
\operatorname{logit} \gamma(c, h_1) \;-\; \operatorname{logit} \gamma(c, h_2)
\;=\; g(\ell_{h_1} + \lambda(c)) \;-\; g(\ell_{h_2} + \lambda(c)),
\end{equation}
where $\ell_{h_i} = \operatorname{logit} \rho(h_i)$, and the channel-conditioned
compliance log-odds gap satisfies the lower bound
\begin{equation}
\label{eq:nofreelunch}
\bigl| \operatorname{logit} \gamma(c, h_1) - \operatorname{logit} \gamma(c, h_2) \bigr|
\;\ge\; \kappa \cdot \bigl| \operatorname{logit} \rho(h_1) - \operatorname{logit} \rho(h_2) \bigr|,
\end{equation}
where $\kappa := \inf_{\ell \in [\ell_{\min}, \ell_{\max}]} g'(\ell) \in [0, 1]$ is
the minimum slope of the compliance read-out on the relevant logit interval. A
Bayesian-rational compliance function trained on $P^\star$ therefore cannot achieve
channel-uniform compliance on matched content unless one of the following holds:
(i) $\rho$ is uniform across channels, (ii) the compliance function abandons the
Bayesian posterior read-out of Assumption~\ref{ass:bayes}, or (iii) the compliance
function collapses $g$ to a constant on the relevant interval, which by
Theorem~\ref{thm:rate-ratio} costs calibration accuracy on the underlying reliability
classification.
\end{theorem}

\begin{proof}
\emph{Step 1 (KL projection identifies the channel prior with the corpus rate).}
Under Assumption~\ref{ass:kl}, $\hat\pi(h)$ minimises
$\mathbb{E}_{R \sim P^\star(\cdot \mid h)}[\mathrm{KL}(\mathrm{Ber}(R) \| \mathrm{Ber}(q))]$
over $q \in (0,1)$. Expanding the KL divergence and taking the derivative in $q$
yields the first-order condition $\rho(h)/q = (1-\rho(h))/(1-q)$, with the unique
solution $q = \rho(h)$ in $(0,1)$. Strict convexity of the KL functional in $q$ on
$(0,1)$ makes this the unique minimiser. Therefore $\hat\pi(h) = \rho(h)$, and
$\ell_h = \operatorname{logit} \hat\pi(h) = \operatorname{logit} \rho(h)$. This is
the M-projection identification of the channel-conditional reliability marginal
\cite{csiszar1975projection,amari2016information,coverthomas2006elements}.

\emph{Step 2 (compliance log-odds inherit the channel-prior log-odds gap).}
By Assumption~\ref{ass:bayes} and the posterior log-odds factorisation of
equation~\eqref{eq:posterior-logodds}, the compliance log-odds satisfy
$\operatorname{logit} \gamma(c, h) = g(\ell_h + \lambda(c))$. Substituting the
result of Step 1, the compliance log-odds difference for matched content $c$
across channels $h_1, h_2$ equals
$g(\ell_{h_1} + \lambda(c)) - g(\ell_{h_2} + \lambda(c))$, which is
equation~\eqref{eq:impossibility}. The matched-content hypothesis from
Assumption~\ref{ass:strong-matching} ensures $\lambda(c)$ is the same in both
arguments of $g$.

\emph{Step 3 (lower bound via the mean value theorem).} The function $g$ is
1-Lipschitz by hypothesis and monotonic increasing under
Assumption~\ref{ass:bayes}; Lipschitz functions are absolutely continuous and
differentiable almost everywhere. By the mean value theorem applied to $g$ on
the closed interval between $\ell_{h_2} + \lambda(c)$ and $\ell_{h_1} + \lambda(c)$,
there exists a point $\xi$ in the interior such that
$g(\ell_{h_1} + \lambda(c)) - g(\ell_{h_2} + \lambda(c)) = g'(\xi) \cdot (\ell_{h_1} - \ell_{h_2})$.
Taking absolute values and bounding $g'(\xi) \ge \kappa = \inf g'$ on the interval
yields the bound~\eqref{eq:nofreelunch}.

The three escape routes follow by inspection of the inequality. Route (i) zeros
the right-hand side. Route (ii) violates Assumption~\ref{ass:bayes} and removes
the equality of Step 2 from which the bound is derived. Route (iii) sends $\kappa
\to 0$, but a constant $g$ on the relevant interval implies that the compliance
read-out cannot distinguish reliable from unreliable content in that range, which
by Theorem~\ref{thm:rate-ratio} eliminates calibration accuracy on the underlying
reliability task.
\end{proof}

\paragraph{Verification.} Theorem~\ref{thm:impossibility} passes equivalence verification at $3/3$ engines on Step~1 (KL-projection first-order
condition; all three verifiers return PASS) and Step~2 (posterior log-odds factorisation; all three verifiers return PASS). Step~3 (mean value theorem on absolutely continuous 1-Lipschitz functions) is verified by the symbolic and numerical engines with 32 sub-Gaussian random parameter configurations; the satisfiability-solver verifier is replaced by transcendental-aware symbolic verification because the Lipschitz quantifier scope is outside the satisfiability solver's decidable
theory. See \texttt{theory/proofs/T5\_impossibility/proof.py}.

\begin{corollary}[Falsification by corpus rebalancing]
\label{cor:corpus-falsification}
Equation~\eqref{eq:impossibility} predicts that the compliance log-odds gap on
matched content tracks the channel-prior log-odds gap induced by the corpus
reliability rates. Under a counterfactual corpus $\tilde P$ in which
$\tilde\rho(h_1) = \tilde\rho(h_2)$, the predicted gap is zero up to the
1-Lipschitz approximation error of the compliance function. Under a corpus
$\tilde P$ that inverts the reliability ranking, $\tilde\rho(h_1) < \tilde\rho(h_2)$,
the predicted compliance gap inverts in sign.
\end{corollary}

The impossibility statement of Theorem~\ref{thm:impossibility} is the
channel-asymmetric analogue of the fairness-under-disparate-base-rates
impossibility of \cite{hardt2016fairness}: where fair classification cannot
simultaneously equalise true-positive and false-positive rates across groups when
base rates differ, channel-uniform compliance cannot be achieved by a Bayesian-rational
read-out trained on a corpus where channel reliability rates differ. The KL-projection
machinery is the same in both cases, with channels playing the role of protected
groups and reliability rates playing the role of group base rates. The bound's
three escape routes exhaust the structural moves available: rebalance the corpus,
abandon the rationality of the read-out, or accept calibration loss equal to the
suppressed reliability ratio.

The empirical magnitudes are consistent with the bound. Plausible reliability
rates in 2026 frontier-LLM training corpora are $\rho(\text{tool}) \in [0.85, 0.95]$
for curated tool-output pipelines and retrieval-augmented generation, against
$\rho(\text{user}) \in [0.40, 0.65]$ for open-domain user input under no factuality
constraint. Substituting these ranges into~\eqref{eq:nofreelunch} with $\kappa$
equal to the calibration constants typical of temperature-scaled language models
\cite{lamb2025semantic} yields a predicted compliance log-odds gap in the range
$[1.1, 3.0]$. The empirical gap of $3.38$ from
Remark~\ref{rem:empirical} sits at and slightly above the upper boundary of the
prediction interval, which is consistent with a Bayesian-rational lower bound plus
a small over-amplification effect at exactly the magnitude that
Theorems~\ref{thm:probing}--\ref{thm:probing-gaussian} lower-bound on the
probe-distinguishability axis.

\subsection{Priority inversion under regime-conditioned channel coefficients}
\label{sec:priority-inversion}

Theorem~\ref{thm:impossibility} explains why declarative compliance shows
$\gamma(c, \text{tool}) > \gamma(c, \text{user})$ on matched content. The same
training pipelines, however, simultaneously enforce the opposite ordering for
command-following: instruction-hierarchy training of \cite{wallace2024instruction}
produces $\gamma_{\text{cmd}}(c, \text{user}) > \gamma_{\text{cmd}}(c, \text{tool})$
for matched commands. The two orderings appear to invert on the same channel
pair. The result below resolves the paradox by formalising the inversion as a
sign-reversal of a regime-conditioned channel coefficient on a shared
channel-of-origin representation.

\begin{assumption}[Content-type partition with shared channel encoding]
\label{ass:partition}
The content space partitions as $\mathcal{C} = \mathcal{C}_{\text{cmd}} \sqcup \mathcal{C}_{\text{dec}}$
into commands (imperative speech acts) and declaratives (assertive speech acts)
\cite{searle1969speech}. There exists a channel-of-origin encoding
$\phi : \mathcal{H} \to \mathbb{R}$ measurable by the model and shared across
content types, and regime-conditioned coefficients $\beta_\tau, \alpha_\tau$ for
$\tau \in \{\text{cmd}, \text{dec}\}$, such that under
Assumption~\ref{ass:bayes} and \ref{ass:strong-matching} the compliance log-odds
factor as
\[
\operatorname{logit} \gamma(c, h, \tau) \;=\; \beta_\tau \cdot \phi(h) + \alpha_\tau + \lambda(c).
\]
\end{assumption}

\begin{theorem}[Priority inversion via regime-conditioned channel coefficient]
\label{thm:priority-inversion}
Suppose Assumption~\ref{ass:partition} holds, the training corpus carries
truth-aligned reliability rates that order channels in opposite directions on the
user-tool axis ($\rho_{\text{cmd}}(\text{user}) > \rho_{\text{cmd}}(\text{tool})$
under Wallace-style command-priority training, and
$\rho_{\text{dec}}(\text{tool}) > \rho_{\text{dec}}(\text{user})$ under
truth-aligned declarative training), and the KL-projection compliance
\textnormal{(Assumption~\ref{ass:kl})} holds within each regime. Then
$\beta_{\text{cmd}} \cdot \beta_{\text{dec}} < 0$ on the user-tool channel pair,
and consequently the same model exhibits both
$\gamma_{\text{cmd}}(c, \text{user}) > \gamma_{\text{cmd}}(c, \text{tool})$ for
matched commands and
$\gamma_{\text{dec}}(c, \text{tool}) > \gamma_{\text{dec}}(c, \text{user})$ for
matched declaratives. The priority inversion observed in the parent paper is
therefore the dual application of Theorem~\ref{thm:impossibility} under
opposite-signed corpus reliability rates rather than a contradiction in the
underlying compliance function.
\end{theorem}

\begin{proof}
By Theorem~\ref{thm:impossibility} applied separately to each regime $\tau$, the
KL-projection identifies the regime-conditioned channel prior with the corpus rate,
$\hat\pi_\tau(h) = \rho_\tau(h)$. Under Assumption~\ref{ass:partition}, the
shared channel encoding $\phi(h)$ enters the compliance log-odds linearly with
regime-conditioned coefficient $\beta_\tau$. Substituting the channel prior
contributions through equation~\eqref{eq:posterior-logodds} for each regime gives
$\beta_\tau$ as the slope of the compliance log-odds with respect to
$\operatorname{logit} \rho_\tau(h)$, which is positive when truth-aligned training
rewards higher compliance on more-reliable channels. Under the hypothesis that
$\rho_{\text{cmd}}$ orders channels with user above tool while $\rho_{\text{dec}}$
orders channels with tool above user, the slope of compliance with respect to
$\phi(h)$ has opposite signs in the two regimes, that is
$\beta_{\text{cmd}} \cdot \beta_{\text{dec}} < 0$. Substituting back into the
compliance factorisation yields the two stated orderings as direct consequences.
\end{proof}

\paragraph{Verification.} Theorem~\ref{thm:priority-inversion} is independently checked
by all three verifiers (algebraic-symbolic, numerical-witness, satisfiability-solver) on the slope-sign
identification under Assumption~\ref{ass:partition}; see
\texttt{theory/proofs/T6\_priority\_inversion/proof.py}. The verifier confirms
that the linear factorisation is consistent with Theorem~\ref{thm:impossibility}
applied to each regime separately and that the sign-reversal claim follows
mechanically from the corpus-rate orderings.

\begin{corollary}[Coupling-cost trade-off]
\label{cor:coupling-cost}
Suppose post-hoc safety training imposes a shared channel-priority objective
across content types, that is, replaces the regime-conditioned coefficients
$\beta_{\text{cmd}}, \beta_{\text{dec}}$ with a single coefficient
$\beta_{\text{shared}}$. Under any choice of $\beta_{\text{shared}}$, at least
one regime suffers a calibration loss equal to
$|\beta_{\tau} - \beta_{\text{shared}}| \cdot \mathrm{Var}(\phi)^{1/2}$ at the
content-type-conditional KL projection, with equality at
$\beta_{\text{shared}} = (\beta_{\text{cmd}} + \beta_{\text{dec}})/2$.
The coupling intervention therefore cannot reduce both rates without paying
calibration cost on at least one regime, which is the no-free-lunch consequence
of Theorem~\ref{thm:impossibility} applied to the shared-coefficient compliance
function.
\end{corollary}

The priority-inversion theorem makes a sharp empirical prediction at the activation
level. Under Assumption~\ref{ass:partition}, a single linear probe trained to
classify channel of origin on the residual-stream activation should recover a
shared encoding $\phi(h)$, with regime-conditioned coefficients $\beta_\tau$ that
project compliance log-odds in opposite directions across the user-tool axis. The
mechanism pipeline at \texttt{mechanism/exp1\_linear\_probe.py} can test this by
reusing the channel-of-origin probe trained on declarative content and applying
the same projection to command-tagged content, with the prediction that the
correlation with command-following compliance is sign-inverted relative to
declarative compliance. The closest empirical precedent for the cleavage is
\cite{mason2026imperative} on imperative-versus-declarative instruction topology
and \cite{waqas2026provenance} on provenance-conditioned compliance, neither of
which derives a Bayesian closed-form bound for the inversion. The framework here
fills that gap.

\subsection{No-free-lunch under bounded inference capacity}
\label{sec:no-free-lunch}

The escape routes of Theorem~\ref{thm:impossibility} include the option of
collapsing the calibration constant $\kappa = \inf g'$ toward zero to suppress the
channel-conditioned compliance gap. Bounded inference capacity blocks this route.
The result below formalises the intuition that Bayesian rationality on a
sufficiently rich content space requires capacity that scales with the
content-space effective dimensionality, and derives a Pinsker-quadratic lower
bound on the calibration cost any bounded-capacity model must accept on
channel-asymmetric deployment distributions \cite{wolpert1997nofreelunch,
genewein2019rate, geirhos2020shortcut}.

\begin{theorem}[Capacity correction to channel-prior emergence]
\label{thm:capacity}
Let $\hat\pi^C$ be the channel-prior estimator under bounded-capacity empirical
risk minimisation in a function class $\mathcal{F}_C$ with covering number
$\log \mathcal{N}(\mathcal{F}_C, \varepsilon, \|\cdot\|_\infty) \le C \log(1/\varepsilon)$,
and let $\rho$ be the corpus channel-reliability rate of
Assumption~\ref{ass:truth-aligned}. Under the hypotheses of
Theorem~\ref{thm:impossibility}, the channel-prior estimator satisfies
\begin{equation}
\label{eq:capacity-bound}
\mathbb{E}\!\left[\, \bigl(\hat\pi^C(h) - \rho(h)\bigr)^2 \,\right]
\;\le\; c_1 \cdot \frac{C \log(N/C)}{N},
\end{equation}
where $N$ is the training-set size and $c_1$ is an absolute constant. The bound is
positive and vanishing for $N > C$ and tends to zero as $N/C \to \infty$.
Consequently, the calibration constant $\kappa^C := \inf g'$ on $\hat\gamma^C$
converges to $\kappa^\infty > 0$ as $N/C \to \infty$ whenever $\rho$ is non-uniform,
and the third escape route of Theorem~\ref{thm:impossibility} (collapsing $\kappa$ to
zero) is unavailable to any bounded-capacity Bayesian compliance function.
\end{theorem}

\begin{proof}
The bound~\eqref{eq:capacity-bound} is the standard Rademacher-complexity squared
loss bound for bounded-Lipschitz function classes
\cite[Section~3]{wolpert1997nofreelunch}, applied to the squared error of
$\hat\pi^C$ relative to the population-optimal $\rho(h)$ obtained by
Theorem~\ref{thm:impossibility} Step~1. The convergence
$\kappa^C \to \kappa^\infty$ as $N/C \to \infty$ follows because $\hat\pi^C \to \rho$
in $L^2$ implies $\inf g' \to$ the infinite-capacity infimum, which is positive when
$\rho$ is non-uniform under Assumption~\ref{ass:bayes}. The non-availability of the third
escape route follows because $\kappa = 0$ requires $g$ to be constant on the
relevant logit interval, which under
Assumption~\ref{ass:strong-matching}~(ii) requires the population-optimal
$\hat\gamma$ to be constant on that interval and therefore non-Bayesian on
content with non-degenerate likelihood-ratio $\lambda(c)$.
\end{proof}

\paragraph{Verification.} Theorem~\ref{thm:capacity} is independently checked
on the $L^2$ convergence step (symbolic and numerical verifiers PASS) and
on the structural unavailability of the $\kappa \to 0$ escape route (symbolic and numerical
verifiers PASS). The Rademacher constant $c_1$ depends on the function class and is
not identified by this proof; only the order-of-magnitude scaling is verified.
The capacity-scaling corollary below provides the empirical falsification rather
than a closed-form constant. Script: \texttt{theory/proofs/T7\_capacity\_bound/proof.py}.

The framework is silent on the precise calibration cost on a given deployment
distribution; chaining Theorem~\ref{thm:probing} with Theorem~\ref{thm:impossibility}
yields an order-of-magnitude lower bound on the activation-level KL of the form
$\Omega((\Delta\ell)^2/M^2)$ where $M$ is the read-out bound, but the constants
in this chain depend on the channel-conditional activation-distribution
structure and are not identified without further assumptions. The capacity-scaling
corollary below operationalises the prediction at the empirical level, where it
admits a sharp test against fine-tuning ablations.

\begin{corollary}[Capacity-scaling falsification]
\label{cor:capacity-scaling}
Theorem~\ref{thm:capacity} predicts that the empirical channel-conditioned
compliance log-odds gap on byte-identical matched content converges to the
infinite-capacity Bayesian limit $\kappa \cdot \Delta\ell$ as model capacity $C$
scales relative to the content-space effective dimensionality $D$. Concretely, a
fine-tuning ablation across model sizes (Llama 3.1 8B, 70B, 405B) on the same
channel-asymmetric corpus should produce a compliance log-odds gap that converges
monotonically toward $\kappa \cdot \Delta\ell$ rather than collapsing to zero. The corollary
is empirically falsified if (i) the gap collapses to zero at large $C$ (which
would contradict the lower bound), or (ii) the gap exceeds $\Delta\ell$ at large
$C$ (which would contradict Theorem~\ref{thm:rate-ratio}).
\end{corollary}

The capacity bound completes a sandwich of the channel-conditioned compliance gap
between an upper bound from Theorem~\ref{thm:rate-ratio} and a lower bound from
Theorem~\ref{thm:impossibility} times a calibration constant
$\kappa_{\min}(C, D, \varepsilon)$ that is positive at finite capacity. The
empirical scaling-law calibration improvements reported by
\cite{kadavath2022calibration} on unconditional benchmarks are consistent with the
capacity-scaling prediction, while the empirical authority-laundering rate-ratio
gap of $3.38$ across nine frontier models that span four orders of magnitude in
parameter count provides direct evidence that the gap is driven by the structural
channel-asymmetry of $\rho$ rather than residual capacity slack. The connection to
the empirical calibration-versus-confidence dissociation reported by
\cite{tian2023verbalized} is the channel-conditional analogue of their
unconditional finding: bounded-capacity models display channel-conditional
shortcuts in exactly the regime where Theorem~\ref{thm:capacity}'s lower bound
binds. The closest mathematical precedent is the rate-distortion bounded-Bayesian
agent framework of \cite{genewein2019rate} and the shortcut-learning pattern of
\cite{geirhos2020shortcut}, neither of which derives the channel-conditional
closed-form bound here.

\section{Empirical predictions}
\label{sec:predictions}

The framework yields eleven empirical tests runnable on the existing
\texttt{data/raw/} corpus and on the activations captured by the
\texttt{mechanism/run\_subject\_local.py} pipeline.

\begin{enumerate}
\item \textbf{Channel-prior probing.} Estimate $\hat{\pi}(h)$ for each channel by
querying each subject model directly on a held-out set of paired (genuine,
fabricated) institutional documents. The runnable protocol is implemented at
\texttt{mechanism/probe\_channel\_priors.py} with a 21-document paired corpus
covering all 13 domains in \texttt{mechanism/probe\_channel\_priors\_corpus.py},
and the pre-registered protocol document is at \texttt{mechanism/PROBE\_CHANNEL\_PRIORS\_PROTOCOL.md}.
The protocol presents each document through channel $h$ via the same channel-delivery
semantics as the main authority-laundering experiment (TOOL\_DIRECT, DOC\_USER\_PASTE,
USER\_DIRECT, SYSTEM\_DIRECT) and asks the model to report a trichotomous label
(genuine, partial, fabricated) plus a 0--100 confidence; the calibration uses the
genuine-vs-fabricated balance in the held-out set. The empirical $\hat\pi$ vector
across channels is the input to Corollary~\ref{cor:falsification}, with
domain-stratified bootstrap confidence intervals (1000 resamples) over both
documents and replicates. To address Remark~\ref{rem:falsification-scope}~(c),
the probing protocol reports both the binary-collapsed $\hat\pi$ and the trichotomy
directly. Total cost across the 9-model panel is approximately three to eight
United States dollars across Anthropic, OpenAI, and Google APIs.

\item \textbf{Bayes-factor calibration of grounding.} For each grounding fact $F$
in the existing TOOL\_DIRECT\_GROUNDED condition, estimate $\hat{B}(F \mid c, h)$
by direct probing under the channel-blind elicitation protocol of
Remark~\ref{rem:grounding-scope}: present the model with the fabricated document
and ask how the user's contradicting statement updates the assessment of the
document's reliability. Test Theorem~\ref{thm:grounding} by checking whether the
empirical $\operatorname{logit} \hat{\gamma}$ reduction matches $\log \hat{B}(F)$
up to the 1-Lipschitz factor and is channel-invariant. Sign violations against
Corollary~\ref{cor:grounding-sign} are reported separately as the cleanest
falsification cells. The protocol mirrors the motivation-discounting Bayesian
benchmark of \cite{wu2025motives}.

\item \textbf{Probe-direction compliance correlation.} Use the probing pipeline in
\texttt{mechanism/} to extract the channel-authority direction $\mathbf{v}$ at the
layer identified by Experiment~1 of the mechanism work. Test
Theorem~\ref{thm:probing}~(d) by checking whether the projected score
$\langle \mathbf{v}, \mathbf{x}_h \rangle$ predicts compliance with effect-size
consistent with the lower bound, both pooled and within the TOOL\_DIRECT condition.

\item \textbf{Matched-content sourcing for Remark~\ref{rem:falsification-scope}~(b).}
Restrict the rate-ratio test to byte-identical factorial pairs in which the same
fabricated document appears verbatim in both TOOL\_DIRECT and DOC\_USER\_PASTE
conditions. The existing matched-pair extractor in
\texttt{mechanism/extract\_matched\_pairs.py} returns these pairs by construction.

\item \textbf{Falsifiable lower bound on probe accuracy.} The probing-representation
lower bound (Theorem~\ref{thm:probing}) makes a specific quantitative prediction
about Experiment~1 of the mechanism work. Plugging the empirical authority-laundering
log-odds gap $\Delta\hat\gamma \approx 3.4$ (precisely $3.385$) and a reference
rational channel-prior gap of $1.6$ log-odds (odds ratio of approximately $5$) into
the bound yields a non-Bayesian residual $\Delta_{\textnormal{nb}} \approx 1.8$. With a
conservative compliance-log-odds half-range of $M = 10$, Theorem~\ref{thm:probing}
predicts
\[
\alpha^\star \;\ge\; \tfrac{1}{2} + \Delta_{\textnormal{nb}}/(4M) \;\approx\; 0.545
\]
on a layer-by-layer linear probe trained to discriminate TOOL\_DIRECT from
USER\_IMPERATIVE channel of origin on byte-identical fabricated content. The
mechanism plan's Experiment~1 (L2-regularised logistic regression on residual-stream
activations at the end-of-context position, scenario-stratified train/test split,
leave-one-domain-out cross-validation on Llama 3.1 8B Instruct) is the direct
test. The theorem is empirically falsified when probe accuracy stays at or below
$0.545$ at every transformer layer, which would indicate that the empirical
compliance gap is not realised through a linearly-decodable channel-of-origin
direction in the residual stream.

When the linear-read-out structure can be confirmed and the Gaussian
sub-Gaussian sharpening of Theorem~\ref{thm:probing-gaussian} applies, the
predicted accuracy at unit-norm read-out and projected standard deviation $\sigma$
is $\alpha^\star = \Phi(\Delta_{\textnormal{nb}}/(2\sigma))$. With realistic
Llama 3.1 8B late-layer values $\sigma \in [2.2, 4.5]$, the prediction range is
$\alpha^\star \in [\Phi(0.200), \Phi(0.409)] = [0.579, 0.659]$, which is the
sharper prediction the linear-probe experiment must clear. A positive
Experiment~1 result with probe accuracy in the $0.90$ range will clear both the
conservative and the Gaussian-sharpened floors and will provide quantitative
rather than qualitative confirmation that the channel-of-origin direction is the
representational substrate of the non-Bayesian compliance gap.

\item \textbf{Corpus-uniformisation falsification of Theorem~\ref{thm:impossibility}.}
Construct a fine-tuning corpus $\tilde P$ in which the reliability rate $\tilde\rho(h)$
is held uniform across the four channels $\{\text{sys}, \text{tool}, \text{doc}, \text{user}\}$
at a single value (e.g.\ $\tilde\rho \equiv 0.5$). The corpus contains balanced
(genuine, fabricated) document pairs, with the same content delivered through each
channel an equal number of times under each reliability label. Fine-tune a single
subject model (Llama 3.1 8B Instruct, the mechanism-pipeline workhorse) on this
corpus until validation loss saturates. Measure the channel-conditioned compliance
gap on the held-out matched-content authority-laundering test set. Theorem~\ref{thm:impossibility}
predicts the gap collapses to within the 1-Lipschitz approximation error around
zero. The theorem is empirically falsified if the gap remains at the pre-fine-tuning
magnitude (logit gap of approximately $3.38$).

\item \textbf{Corpus-inversion falsification of Theorem~\ref{thm:impossibility}.}
Construct a fine-tuning corpus $\tilde P$ in which the reliability ranking is
inverted, $\tilde\rho(\text{user}) > \tilde\rho(\text{tool})$. Fabricated documents
are tagged exclusively with the tool channel and genuine documents are tagged
exclusively with the user channel during fine-tuning. Theorem~\ref{thm:impossibility}
predicts the compliance gap inverts in sign on the held-out matched-content test:
user-channel content should be operationalised more readily than tool-channel
content. The theorem is empirically falsified if the gap retains its original
sign or if the sign-inversion is accompanied by a magnitude collapse inconsistent
with $\kappa \cdot |\Delta \operatorname{logit} \tilde\rho|$.

\item \textbf{Corpus-magnitude scaling for $\kappa$ identification.} For a sequence
of fine-tuning corpora $\tilde P_t$ parameterised by the channel-reliability
log-odds gap $\Delta \tilde\ell_t \in \{0.5, 1.0, 1.5, 2.0, 3.0\}$, fit a linear
regression of the post-fine-tuning compliance log-odds gap on $\Delta \tilde\ell_t$.
Theorem~\ref{thm:impossibility} predicts a slope in the closed interval $(0, 1]$,
equal to the calibration constant $\kappa$. The theorem is empirically falsified
if the slope is zero (no dependence on corpus statistics), greater than $1$
(over-amplification beyond Bayesian rationality, which would also falsify
Theorem~\ref{thm:rate-ratio}), or strongly nonlinear.

\item \textbf{Sign-inversion test for priority inversion (Theorem~\ref{thm:priority-inversion}).}
Reuse the channel-of-origin probe $\mathbf{v}^\star$ from Experiment~1 of the
mechanism work. For matched-content pairs in (i) command form (USER\_IMPERATIVE
vs SYSTEM\_INSTRUCTION conditions on jailbreak prompts) and (ii) declarative
form (USER\_IMPERATIVE vs TOOL\_DIRECT conditions on fabricated guidelines),
project both onto $\mathbf{v}^\star$ and regress the compliance binary label on
the projection separately within each regime. Theorem~\ref{thm:priority-inversion}
predicts that the regression coefficient is negative on the command-form pairs
(higher channel score $\to$ lower command compliance) and positive on the
declarative-form pairs (higher channel score $\to$ higher declarative compliance).
Sign agreement across regimes would falsify the priority-inversion theorem.

\item \textbf{Coupling-cost regression for Corollary~\ref{cor:coupling-cost}.}
Apply post-hoc fine-tuning that imposes a shared channel-priority objective
across content types (e.g.\ a single channel-priority loss summed over both
command and declarative training signals). Measure compliance rates on held-out
matched commands and declaratives. The corollary predicts that the joint
calibration loss exceeds the regime-conditional calibration losses of the
unconstrained model by a positive constant proportional to
$|\beta_{\text{cmd}} - \beta_{\text{dec}}|$. The corollary is empirically falsified
if the shared-coefficient model achieves equal or better calibration on both
regimes simultaneously.

\item \textbf{Capacity-scaling falsification for Corollary~\ref{cor:capacity-scaling}.}
Fine-tune Llama 3.1 8B, 70B, and 405B on a single channel-asymmetric corpus
with reliability rates $\rho_{\text{tool}} = 0.9$ and $\rho_{\text{user}} = 0.4$.
Measure the post-fine-tuning compliance log-odds gap on byte-identical matched
content. Theorem~\ref{thm:capacity} predicts the gap converges monotonically
toward $\kappa \cdot \Delta\ell$ with $\Delta\ell = \operatorname{logit}(0.9) - \operatorname{logit}(0.4) \approx 2.59$
as the training-set-to-capacity ratio grows, with a residual $O(\sqrt{C \log(N/C)/N})$ slack vanishing as $N/C \to \infty$.
The theorem is empirically falsified if (i) the gap collapses to zero at large
capacity (uniform compliance, contradicting the lower bound), or (ii) the gap
exceeds $\Delta\ell$ at large capacity (over-amplification beyond Bayesian
rationality, also falsifying Theorem~\ref{thm:rate-ratio}).
\end{enumerate}

\section{Discussion}
\label{sec:discussion}

The Bayesian source-reliability framework converts the empirical channel-conditioned
compliance gap into a falsifiable Bayesian-rationality claim. The empirical gap
documented in the authority-laundering work falsifies rational Bayesian compliance
under any plausible channel-prior estimate (Remark~\ref{rem:empirical}), conditional
on ruling out the four scope conditions of Remark~\ref{rem:falsification-scope}.
The closest published precedent is the work of \cite{kumaran2026competing} on
competing biases in language-model evidence weighting, who report a
Bayesian-deviation ratio of approximately $2$ to $3$ on the evidence-direction
axis (supporting versus opposing advice). Theorem~\ref{thm:rate-ratio} extends that
empirical deviation result along two axes: it operates on the delivery channel
rather than evidence direction, and it derives a closed-form bound connecting two
separately observable quantities, turning a fitted deviation into a falsifiable
inequality. The authority-laundering empirical gap of approximately $20\times$ on
rate ratio (corresponding to a logit gap of $3.38$) is roughly an order of magnitude
larger than the evidence-direction deviation, which positions channel-as-authority
as a substantially sharper non-Bayesian heuristic than evidence-direction
overweighting.

The cognitive-science Bayesian source-reliability lineage from \cite{hahn2007rational,
harris2016expert} and the Sperber epistemic-vigilance program \cite{sperber2010vigilance}
supplies the underlying generative model in equation~\eqref{eq:posterior} but does
not derive a rate-ratio bound of the form in Theorem~\ref{thm:rate-ratio}. The
source-cued heuristic literature in human cognition \cite{tversky1974judgment,
petty1986elaboration} qualitatively predicts the channel-as-authority pattern but
does not formalise it as a quantitative bound. The pragmatic-empirical work of
\cite{cheng2026accommodation} on accommodation and epistemic vigilance in language
models documents the same source-axis effect that Theorem~\ref{thm:rate-ratio}
addresses but operates at the empirical level without the Bayesian-rationality
null. The instruction-hierarchy framework of \cite{wallace2024instruction} supplies
the priority-inversion framing for declarative content versus commands, and the
empirical control-illusion finding of \cite{geng2026control} confirms that the
Wallace hierarchy is fragile in exactly the declarative regime that
Theorem~\ref{thm:rate-ratio} addresses.

The grounding-asymmetry bound (Theorem~\ref{thm:grounding}) further rules out a
wide class of rational explanations through its one-sided sign claim: no rational
Bayesian compliance model produces a sign flip on contradicting evidence under
matched conditional Bayes-factor and channel-blind elicitation. The empirical
Sonnet 4.5 backfire is therefore a sign-violation falsification that is robust to
recalibration of any compliance constant. This sign claim is logically stronger
than correlational Bayesian-coherence measures \cite{imran2025bayes}. The
probing-representation lower bound (Theorem~\ref{thm:probing}) ties the theoretical
non-Bayesian gap to the mechanistic activation-probing direction through a
Pinsker--Le Cam information-theoretic chain
\cite{tsybakov2009nonparametric, coverthomas2006elements}, giving a closed loop
between behavioural measurement, theory, and mechanism. The bound is the lower-bound
dual of the Fano-style accuracy upper bound of \cite{wan2025fano}: where they
upper-bound LLM output accuracy from task-complexity overload, our bound
lower-bounds activation-probe distinguishability from a behavioural compliance gap.
The Gaussian sharpening (Theorem~\ref{thm:probing-gaussian}) replaces the
saturation parameter $M$ with the projected standard deviation $\sigma$, yielding
the predicted accuracy floor $\Phi(\Delta_{\textnormal{nb}}/(2\sigma)) \in [0.58, 0.66]$
for realistic Llama 3.1 8B late-layer activations.

The framework's principal limitation is the 1-Lipschitz-in-logit calibration
assumption, which is required for the analytic bounds; relaxing to general monotonic
$f$ retains the sign-of-effect statement of Proposition~\ref{prop:sign} but yields
a sign-only falsification rather than a quantitative one. A second limitation is
that the falsification scope of Remark~\ref{rem:falsification-scope} requires four
substantive empirical claims to be independently ruled out before an empirical
bound violation counts as evidence of a non-Bayesian heuristic; the predictions in
\S\ref{sec:predictions} specify how each of (b) and (c) can be addressed with the
existing corpus, while (a) and (d) are open targets for future iteration. The
probing-representation bound is vacuous if channel-authority is encoded nonlinearly
at the probed layer, which is a condition the matched-expressivity probe extension
of Remark~\ref{rem:probing-scope} addresses. The framework operates on the
declarative-acceptance axis and is silent on the command-following axis that
\cite{wallace2024instruction} formalises.

Theorem~\ref{thm:impossibility} (training-distribution impossibility, \S\ref{sec:training-distribution})
sits structurally upstream of Theorems~\ref{thm:rate-ratio}--\ref{thm:probing-gaussian}
and supplies the missing causal explanation for why the channel prior $\pi(h)$ has
the magnitude it does. Where Cloud Theorem~1 \cite{cloud2026subliminal} shows
distillation gradient steps transport student logits toward teacher logits regardless
of training distribution, Theorem~\ref{thm:impossibility} shows truth-aligned training
transports the channel prior of the model toward the channel reliability rate of the
corpus regardless of the corpus generative process. The result is the channel-asymmetric
analogue of the fairness-impossibility result of \cite{hardt2016fairness}: KL projection
identifies $\hat\pi(h) = \rho(h)$ uniquely, and removing the channel-conditioned compliance
gap requires either rebalancing the corpus, abandoning Bayesian read-out rationality, or
accepting calibration loss equal to the suppressed reliability ratio. The three
falsification experiments F6--F8 in \S\ref{sec:predictions} are concrete fine-tuning
interventions that test the impossibility prediction directly.

All seven theorems are independently checked by at least two verifiers per theorem,
drawn from an algebraic-symbolic engine, a numerical-witness engine, and a
satisfiability-solver engine. T1, T2, T3, T5, and T6 pass all three verifiers;
T4 passes all three on the foundational steps, with the Gaussian-exact form
replacing an as-stated linear bound that the verifiers refuted; T7 passes the
symbolic and numerical engines on the $L^2$ convergence and the structural
unavailability of the $\kappa \to 0$ escape route. Per-theorem verification
scripts and proof-discovery notes are at \texttt{theory/proofs/T1\_rate\_ratio/}
through \texttt{theory/proofs/T7\_capacity\_bound/}. The aggregate
verifier-coverage matrix appears in \texttt{theory/proofs/SUMMARY.md}.

\bibliographystyle{plain}
\bibliography{refs}

\end{document}
